\chapter{Удиртгал}
\label{ch:introduction}

%======================================================================
% 2026.11.11 (3): бүтэц, гарчгийг сургуулийн өмнөх дипломтой (Asura-lab/diplom)
% ижил болгосон: 1.1 Үндэслэл, ач холбогдол; 1.2 Зорилго, зорилт (дэд хэсэгтэй);
% 1.3 Судалгааны асуулт ба хэмжигдэх шалгуур. Үг хэллэгийг Монгол 7 дипломын
% удиртгалтай тулгасан (papers/28-mn-thesis-style-20261111.mn.md). Агуулга хуулаагүй.
% 2026.11.11 (2): үндэслэлийг баримтаар, зорилтыг судалгааны ажлын хэлбэрээр бичсэн;
% дуу таних жижиг хэмжилтийг зорилт, судалгааны асуултаас хассан (3.6, 4.1-д үлдсэн).
% Өмнөх бүтэн хувилбар: thesis/archive-20261111/.
%======================================================================

\section{Үндэслэл, ач холбогдол}
\label{sec:intro-background}

Шинэ Монгол Технологийн Коллежийн тухай мэдээлэл нэг газарт төвлөрөөгүй. Элсэлтийн журам, сургалтын төлбөр, хөтөлбөрийн танилцуулга сургуулийн вэб сайтад бий. Харин аль өрөө хэдэн давхарт байдаг, аль анги хэдэн цагт хаана хичээллэдэг, ямар асуудлаар аль албаны хэнд хандах вэ гэдэг нь сайтад байхгүй. Эдгээр нь сургалтын албаны өрөөний жагсаалт, хичээлийн хуваарь, ажилтны жагсаалтад хадгалагддаг. Иймээс сургуульд ирсэн оюутан, элсэгч, эцэг эх, зочин ийм мэдээллийг ажилтнаас биечлэн асууж авдаг.
% ӨӨРИЙН ҮГЭЭР: ресепшнээс цуглуулсан бодит асуултын 2–3 жишээ, аль төрлийн асуулт хамгийн
% их давтагддагийг энд нэм (боломжтой бол хэдэн хүнээс, хэдийд асуусныг хүснэгтээр).
% Агент асуулт, тоо зохиохгүй.

Тулгамдаж буй асуудлыг тодруулахын тулд 2026 оны 11-р сард сургалтын албаны ажилтнуудтай уулзав. Ажилтнууд хоёр зүйлийг онцолсон. Нэгдүгээрт, тэд ижил мэдээллийг өөр өөр хүнд дахин дахин тайлбарладаг. Хоёрдугаарт, нэгдүгээр давхарт цаасан мэдээлэл байрлуулсан боловч хүмүүс түүнээс хариултаа олохгүй, ажилтнаас дэлгэрүүлж асуудаг. Мөн ийм систем зөвхөн ресепшнд биш, гар утаснаас нээгддэг байвал илүү олон хүнд хүрнэ гэж зөвлөсөн.
% ӨӨРИЙН ҮГЭЭР: уулзсан өдөр, хэнтэй (албан тушаал) уулзсаныг нэм.

Ийм бэрхшээл бусад сургуульд ч ажиглагдсан. Тухайлбал, Милано-Бикокка их сургуулийн оюутнууд анги танхим, хичээлийн хуваарь, шалгалтын өрөөний мэдээллийг олоход хэцүү гэж хариулсан \cite{antico2024unimib}. Burch их сургуулийн оюутнууд мэдээллээ сайтаас олж чадаагүй үедээ энгийн асуултаар захиргааны ажилтанд ханддаг бөгөөд энэ нь ажилтны ачааллыг нэмдэг \cite{matic2026burch}. Эдгээр сургууль асуудлыг оюутны асуултад сургуулийн баримтаас хайж хариулдаг туслах системээр шийдэхийг зорьсон.

Ийм туслахыг Монгол хэл дээр хийхэд хоёр бэрхшээл тулгардаг. Нэгдүгээрт, Монгол хэл яриа танилтын нөөц багатай хэлэнд тооцогддог. Олон хэлний Whisper загвар Монгол яриаг 100\%-иас их үгийн алдаатай таньсан \cite{radford2023whisper}, Халх Монгол хэлний нээлттэй сургалтын яриа хэдхэн цаг л бий \cite{mengke2025khalkha}. Хоёрдугаарт, сургуулийн нэрийн өмнөөс өгөх хариулт үнэн зөв байх ёстой. Хэлний загвар байхгүй баримтыг зохиож болдог тул хариултыг баримтын сангаас хайж, олсон эхдээ тулгуурлан өгөх шаардлагатай \cite{lewis2020rag}.

Иймд сургуулийн тархай бутархай мэдээллийг нэг бүтцэд оруулж, Монгол хэлээр асуусан асуултад түүнээс хариулдаг, хариулт байхгүй үед хэнд, хаана хандахыг зааж өгдөг туслах систем боловсруулах шаардлага тодорхой болов.

Энэ ажлын практик ач холбогдол нь сургуулийн ажилтнуудын ижил мэдээллийг давтан тайлбарлах ачааллыг бууруулж, сургуульд ирсэн хүн хэрэгтэй мэдээллээ нэг газраас авах боломжийг бүрдүүлэхэд оршино. Судалгааны ач холбогдол нь гараар зохион байгуулсан мэдлэгийн модыг хайлтад ашиглах аргыг Монгол хэл дээр, сургуулийн бодит мэдээллээр шалгахад оршино.

\section{Зорилго, зорилт}
\label{sec:intro-goals}

\subsection{Судалгааны зорилго}

Энэ судалгааны зорилго нь сургуулийн ресепшнд ирсэн хүний Монгол хэлээр ярьж эсвэл бичиж асуусан асуултад \textbf{сургуулийн албан мэдээлэлд тулгуурлан хариулдаг туслах системийг боловсруулахад} оршино.

\subsection{Судалгааны зорилтууд}

Энэ зорилгод хүрэхийн тулд дөрвөн зорилтыг дэвшүүлэв:
\begin{enumerate}
    \item Ресепшнд ирдэг хүмүүсийн мэдээллийн хэрэгцээ, системд тавигдах шаардлагыг сургуулийн албадтай уулзаж тодорхойлох.
    \item Сургуулийн мэдээллийг хоёр эхээс, тухайлбал албан вэб сайт болон сургуулийн албадаас цуглуулж, мэдлэгийн мод хэлбэрээр зохион байгуулах.
    \item Мэдлэгийн модонд тулгуурласан хайлтын аргыг боловсруулж, бодит асуултаар үнэлэх.
    \item Дуугаар болон бичгээр асуусан асуултад эх сурвалжтай хариулт өгөх системийг хэрэгжүүлэх.
\end{enumerate}
% Зорилт 1: карьерын төв, ресепшний ажилтантай уулзсаны дараа «сургуулийн албад»-ыг нэрээр нь бич.

\section{Судалгааны асуулт ба хэмжигдэх шалгуур}
\label{sec:rq}

Гуравдугаар зорилтыг хэмжихүйц шалгууртай холбохын тулд дараах судалгааны асуултыг тодорхойлов: асуултын сэдвийг мэдлэгийн модны хүрээгээр тодорхойлж, тухайн хүрээний эхэд давуу эрх өгөх нь бүх эхээс ялгалгүй хайхтай харьцуулахад зөв эхийг олох чадварыг сайжруулах уу?

Хэмжих шалгуур нь зөв эх хайлтын эхний нэг болон эхний гурван үр дүнд орсон асуултын тоо юм. Хоёр аргыг ижил асуултаар харьцуулна. Туршилтын аргыг \ref{sec:exp-tree}-т тодорхойлов.

\section{Судалгааны хамрах хүрээ}
\label{sec:intro-scope}

Систем зөвхөн сургуулийн албан мэдээллээр хариулна. Мэдээллийг хоёр эхээс авсан: сургуулийн вэб сайтын нийтэд нээлттэй хуудсууд болон сургуулийн албадаас өгсөн өрөө, анги бүлэг, ажилтан, хичээлийн хуваарийн жагсаалт. Асуулт, хариулт Монгол хэлээр байна. Дүн, ирц зэрэг хувь хүний мэдээлэл болон ажилтны хувийн утасны дугаарыг энэ ажилд хамруулаагүй.

\section{Судалгааны ажлын бүтэц}
\label{sec:intro-structure}

Энэ ажил таван бүлгээс бүрдэнэ. \ref{ch:literature}-р бүлэгт Монгол хэлний яриа танилт, хайлтаар баяжуулсан хариулт, хариулт байхгүй асуулт, мэдлэгийн мод дээрх хайлтын талаарх өмнөх судалгааг тоймлов. \ref{ch:method}-р бүлэгт системд тавигдах шаардлага, мэдээлэл цуглуулсан арга, системийн бүтэц, туршилтын аргыг тодорхойлов. \ref{ch:results}-р бүлэгт одоогоор гарсан үр дүнг, \ref{ch:conclusion}-р бүлэгт дүгнэлт, хязгаарлалт, цаашдын ажлыг тусгав.
